% !TeX program = lualatex
% !TeX encoding = utf8
% !TeX spellcheck = uk_UA
% !TeX root =../PyscfBook.tex


%% --------------------------------------------------------
\section{Проблеми збіжності SCF та їх усунення}\label{sec:convergence}
%% --------------------------------------------------------

Розрахунки у квантовій хімії не завжди проходять гладко. Навіть для простих систем ітераційна процедура самозгодженого поля (SCF) може не досягати стабільного рішення.
Найчастіше проблеми збіжності виникають для перехідних металів, відкритих оболонок та систем із виродженими орбіталями.
У цьому розділі розглянемо основні джерела труднощів та методи їх подолання.

\subsection{Причини поганої збіжності}

Типові джерела нестабільності SCF:
\begin{itemize}
  \item сильна електронна кореляція у незаповнених $d$- і $f$-оболонках;
  \item мала різниця енергій між орбіталями (виродження чи квазівиродження);
  \item невдалий початковий вектор коефіцієнтів (погане стартове наближення);
  \item занадто жорсткі або занадто м'які критерії збіжності.
\end{itemize}

У таких випадках енергія може осцилювати або навіть зростати замість того, щоб збігатися. Для вирішення цих проблем існує кілька стратегій стабілізації.

\subsection{Методи стабілізації збіжності}

Коли стандартний SCF з DIIS-процедурою (див. додаток~\ref{sec:DIIS}) не працює, у PySCF доступні наступні інструменти:

\subsubsection{Level shifting (зсув рівнів)}

Метод штучно підвищує енергії віртуальних орбіталей на величину $\Delta$:
\begin{equation}
\varepsilon_{\text{virt}} \to \varepsilon_{\text{virt}} + \Delta
\end{equation}

Це запобігає коливанням заповнення орбіталей між ітераціями та стабілізує збіжність. Фізично це еквівалентно збільшенню енергетичного бар'єру для переходу електронів з заповнених на віртуальні орбіталі.

\begin{minted}{python}
mf.level_shift = 0.5  # типові значення 0.2-1.0 Hartree
\end{minted}

\subsubsection{Налаштування та альтернативи DIIS у \texttt{PySCF}}

Для прискорення збіжності \texttt{SCF}-процедури у \texttt{PySCF} реалізовано кілька варіантів \texttt{DIIS}-алгоритмів.
Типовим є класичний \texttt{CDIIS}, однак у складних випадках корисно застосовувати розширені схеми.

\paragraph{EDIIS (Energy DIIS)}
Мінімізує енергію замість похибки самоузгодженості. Це часто стабілізує збіжність на ранніх ітераціях, коли система ще далека від стаціонарного стану:
\begin{minted}{python}
mf.DIIS = scf.EDIIS
\end{minted}

\paragraph{ADIIS (Augmented DIIS)}
Комбінує енергетичний (\texttt{EDIIS}) та звичайний (\texttt{DIIS}) підходи.
На початкових ітераціях використовується \texttt{EDIIS}, а ближче до мінімуму автоматично переходить у стандартний \texttt{DIIS}:
\begin{minted}{python}
mf.DIIS = scf.ADIIS
\end{minted}

\paragraph{Параметри DIIS}
Можна тонко контролювати простір інтерполяції, початок застосування DIIS, а також файл для збереження стану:
\begin{minted}{python}
mf.diis_space = 8           # розмір підпростору DIIS
mf.diis_start_cycle = 4     # номер ітерації, з якої почати DIIS
mf.diis_file = 'o2_diis.h5' # файл для збереження стану
\end{minted}

\paragraph{Приклад з офіційного репозиторію \texttt{PySCF}}
Повне налаштування параметрів \texttt{DIIS} наведено у демонстраційному файлі
\inputcode[base=https://github.com/pyscf/pyscf/blob/master/examples/scf/]{24-tune_diis.py}.
У цьому прикладі показано, як переключатися між \texttt{CDIIS}, \texttt{EDIIS} та \texttt{ADIIS},
а також як зберігати та відновлювати стан ітерацій:
\begin{minted}{python}
my_diis_obj = scf.ADIIS()
my_diis_obj.space = 12
my_diis_obj.filename = 'o2_diis.h5'
mf.diis = my_diis_obj
mf.run()
\end{minted}


%\subsubsection{Альтернативні методи екстраполяції}
%
%\paragraph{EDIIS (Energy DIIS)} мінімізує енергію замість похибки самоузгодженості. Це краще працює на початкових ітераціях, коли система далека від збіжності:
%
%\begin{minted}{python}
%mf.DIIS = scf.EDIIS
%\end{minted}
%
%\paragraph{ADIIS (Augmented DIIS)} автоматично комбінує DIIS та EDIIS залежно від стану збіжності --- використовує EDIIS на початку та переключається на DIIS ближче до мінімуму:
%
%\begin{minted}{python}
%mf.DIIS = scf.ADIIS
%\end{minted}
%
%\subsubsection{Налаштування DIIS}
%
%Іноді корисно змінити параметри самого DIIS:
%
%\begin{minted}{python}
%mf.diis_space = 4          # менший підпростір (4 замість 6-8)
%mf.diis_start_cycle = 5    # пізніший старт DIIS
%mf.diis = None             # повністю вимкнути DIIS
%\end{minted}
%
%Зменшення розміру DIIS-простору обмежує «пам'ять» алгоритму, що може допомогти у випадках, коли історія попередніх ітерацій вводить в оману.
%
%\subsubsection{Краще початкове наближення}
%
%Якість початкового наближення критично важлива для збіжності. PySCF пропонує кілька варіантів:
%
%\begin{minted}{python}
%# Стандартні методи
%mf = scf.UHF(mol)
%mf.init_guess = 'minao'  # мінімальний АО (за замовчуванням)
%mf.init_guess = 'atom'   # атомні густини (краще для металів)
%mf.init_guess = '1e'     # лише одноелектронний гамільтоніан
%
%# Використання попереднього розрахунку
%dm_init = mf_previous.make_rdm1()
%mf.kernel(dm0=dm_init)
%\end{minted}
%
%Варіант \inlinecode{'atom'} часто найкращий для перехідних металів та систем з відкритими оболонками.

\subsubsection{Метод Ньютона--Рафсона}

Для систем, які вже близькі до збіжності, але повільно сходяться, можна використати метод другого порядку:

\begin{minted}{python}
mf = mf.newton()  # перехід до Newton-Raphson
mf.kernel()
\end{minted}

Це забезпечує квадратичну збіжність поблизу мінімуму, але вимагає додаткових обчислень гессіану.

\subsection{Вироджені орбіталі та дробові заповнення}

Якщо в системі наявні вироджені або майже вироджені орбіталі, SCF може «коливатись», не вибираючи між ними. Класичні приклади --- атоми з частково заповненими $p$- або $d$-оболонками.

Для таких випадків ефективним є застосування \textit{дробових заповнень орбіталей} (fractional occupations), що згладжують різницю між рівнями енергії.

Ідея полягає у введенні ефективного теплового розмазування Фермі--Дірака:
\[
f_i = \frac{1}{1 + e^{(\varepsilon_i - \mu)/kT}},
\]
де $f_i$ — часткове заповнення орбіталі $i$, $\varepsilon_i$ — її енергія, $\mu$ — хімічний потенціал, а $kT$ — параметр розмазування (зазвичай $\sim 0.1$--$0.3$ Hartree).

Цей підхід дозволяє SCF уникнути нестійких перестрибувань між виродженими рівнями, імітуючи систему при скінченній температурі.

У PySCF для цього достатньо викликати:
\begin{minted}{python}
from pyscf import scf
mf = scf.addons.frac_occ(mf)
\end{minted}

Методика (\inputcode{uhf_fractional_occupations.py}) добре працює для атомів (\ce{V}, \ce{Cr}), радикалів та малих кластерів перехідних металів.

\subsection{Перехідні метали --- приклад складної системи}

Перехідні елементи (Fe, Co, Ni, Cr, Mn тощо) --- класичні приклади систем із поганою збіжністю SCF. Вони поєднують декілька факторів складності одночасно.

\begin{minted}{python}
from pyscf import gto, scf

mol = gto.M(atom="Ni 0 0 0", basis="sto-3g", spin=0)
mf = scf.UHF(mol).run()
print("Converged?", mf.converged)  # часто False
\end{minted}

Причини проблем:
\begin{itemize}
  \item близькість енергетичних рівнів $3d$ і $4s$-орбіталей;
  \item можливість декількох спінових станів, близьких за енергією;
  \item наявність кількох локальних мінімумів функціоналу енергії;
  \item сильна кореляція в незаповненій $d$-оболонці.
\end{itemize}

У PySCF це проявляється як:
\begin{itemize}
  \item осциляції енергії між ітераціями;
  \item розбігання DIIS;
  \item стрибки значення $\langle S^2 \rangle$ між кроками (спін-забруднення).
\end{itemize}

В файлі \inputcode{Transition_Metals_UHF.py} наведено приклад універсальної функції для стабільного розрахунку атомів перехідних металів із урахуванням описаних вище прийомів:

\begin{itemize}
  \item використовується UHF (для врахування спіну);
  \item застосовується зсув рівнів (\inlinecode{level\_shift});
  \item розширюється DIIS-простір (\inlinecode{diis\_space});
  \item використовується атомне початкове наближення (\inlinecode{init\_guess='atom'});
  \item у разі потреби вмикається уточнення Ньютона--Рафсона;
  \item контролюється спін-забруднення через $\langle S^2 \rangle$.
\end{itemize}

Такий комплексний підхід забезпечує стабільність навіть для важких атомів, де стандартний SCF часто не сходиться.

\subsection{Загальна стратегія досягнення збіжності}

Для надійного отримання фізично коректного рішення рекомендується послідовна стратегія --- від простих прийомів до складніших:

\begin{enumerate}
  \item Спробувати стандартний UHF/RHF;
  \item Якщо не збігається --- додати зсув рівнів: \inlinecode{level\_shift=0.5};
  \item Змінити початкове наближення: \inlinecode{init\_guess='atom'};
  \item Використати ADIIS: \inlinecode{mf.DIIS = scf.ADIIS};
  \item Збільшити DIIS-простір: \inlinecode{diis\_space=10};
  \item Для вироджених систем --- дробові заповнення: \inlinecode{frac\_occ};
  \item Якщо близько до збіжності --- метод Ньютона: \inlinecode{mf.newton()};
  \item Вимкнути симетрію: \inlinecode{symmetry=False} (як останній засіб).
\end{enumerate}

В програмі \inputcode{uhf_convergence_strategies.py} продемонстровано різні стратегії покращення збіжності SCF на конкретних прикладах.

\paragraph{Практичні рекомендації:}
\begin{itemize}
  \item Завжди перевіряйте \inlinecode{mf.converged} після розрахунку;
  \item Для систем з відкритими оболонками контролюйте $\langle S^2 \rangle$ --- велике спін-забруднення вказує на проблеми;
  \item Для великих систем на перших етапах можна зменшити точність:\\ \inlinecode{conv\_tol=1e-5};
  \item Якщо енергія осцилює --- це сигнал використати level shifting або EDIIS;
  \item Не бійтесь експериментувати з комбінаціями методів --- іноді неочевидна комбінація працює найкраще.
\end{itemize}

\noindent Таким чином, навіть якщо стандартна процедура SCF не збігається, послідовне застосування описаних прийомів практично гарантує стабільне та фізично обґрунтоване рішення. Ключ до успіху --- розуміння природи проблеми та систематичний підхід до її вирішення.